\section{Proposed Method}

\subsection{Overview and motivation}

The grey wolf optimizer derives its entire search direction from three positions:
the alpha, beta and delta wolves. Every other wolf moves toward the triangle
these three define. Canonical GWO nevertheless spends one true evaluation per
wolf per generation, i.e. $N$ evaluations to maintain a hierarchy that only
three wolves are allowed to influence. In expensive black-box optimization the
opposite asymmetry holds: objective evaluations dominate the cost by orders of
magnitude, whereas generating a candidate position is a few floating point
operations.

This paper therefore asks a question that is orthogonal to the design of new
GWO search operators:

\begin{quote}
\emph{Given a fixed budget of expensive true evaluations, which of the many
candidates a GWO can generate should receive them?}
\end{quote}

We answer it with the Resource-Aware Sparse-Evaluation Grey Wolf Optimizer
(RASE-GWO). Two principles drive the design.

\emph{Decoupling.} The search state of a wolf and its verified state are stored
separately. A wolf's \emph{virtual} position $Y_i$ moves every round and costs
nothing; its \emph{verified} position $X_i^{v}$ is the most recent position that
has actually been confirmed by the objective function. Only verified positions
may compete for the leader hierarchy. This is what makes sparse evaluation
sound for GWO: a leader-driven metaheuristic has no personal-best memory, so
naively reusing unevaluated positions as leaders would let unverified movement
contaminate the search direction. Personal-best based algorithms
(PSO, CSO) are protected by their own memory; GWO is not.

\emph{Over-sampling.} Because a virtual wolf is cheap, we maintain a pool of
$M = NK$ virtual wolves but spend only $K$ true evaluations per round. The
evaluation quota is therefore assigned to a candidate pool that is
$\kappa = M/K$ times larger than the population whose ranking it has to
establish. The extra cost is arithmetic; the budget of expensive evaluations is
unchanged.

\subsection{Dual-state wolves}

Let $M$ be the number of virtual wolves, $D$ the problem dimension and
$[\mathrm{lb},\mathrm{ub}]$ the box. RASE-GWO maintains
\begin{itemize}
  \item $Y \in \R^{M\times D}$: virtual (search) positions;
  \item $X^{v} \in \R^{M\times D}$: verified positions;
  \item $f^{v} \in \R^{M}$: true fitness of each verified position.
\end{itemize}
Only the first $K$ wolves are evaluated during initialization, so the
initial budget, the initial leader triangle and the initial point set are
identical to those of the baseline it is compared against (the harness draws one
shared point pool per problem instance and every algorithm consumes a prefix of
it).

\subsection{Round structure}

Each round performs three steps.

\paragraph{(a) Virtual movement (no evaluations).}
With $a = 2-2\,\FE/\mathrm{MaxFE}$ and independent per-wolf, per-dimension
uniform draws $r_1, r_2 \sim \mathcal{U}(0,1)^D$,
\begin{equation}
\begin{aligned}
A_1 &= 2 a\, r_1 - a, \quad C_1 = 2 r_2, \\
Y^{(1)} &= X_\alpha - A_1 \odot \lvert C_1 \odot X_\alpha - Y \rvert,
\end{aligned}
\end{equation}
and analogously $Y^{(2)}$, $Y^{(3)}$ from $\beta$ and $\delta$. The virtual
population is then
\begin{equation}
Y \leftarrow \clip\!\left(\tfrac{1}{3}\bigl(Y^{(1)}+Y^{(2)}+Y^{(3)}\bigr)\right).
\label{eq:move}
\end{equation}
Equation~\eqref{eq:move} is the published GWO movement; no term is added.

\paragraph{(b) Geometric screening.}
The leader consensus centre is
\begin{equation}
X_{L} = \tfrac{1}{3}\bigl(X_\alpha + X_\beta + X_\delta\bigr).
\end{equation}
Two scores are computed for every virtual wolf:
\begin{align}
d^{c}_{i} &= \lVert Y_i - X_{L} \rVert_2, \label{eq:dc}\\
d^{m}_{i} &= \lVert Y_i - X_{\mu} \rVert_2,
\quad X_{\mu} = \argmax_{X \in \{X_\alpha, X_\beta, X_\delta\}} \lVert X - X_L \rVert_2 .
\label{eq:dm}
\end{align}
$K_c = \lfloor s K \rceil$ slots are given to the $K_c$ wolves with the smallest
$d^{c}$ (\emph{consensus refinement}) and the remaining $K_d = K-K_c$ slots to
the $K_d$ wolves closest to the \emph{minority leader} $X_{\mu}$
(\emph{minority-leader probing}). The evaluation set is
$E = \mathrm{unique}(E_c \cup E_m)$, topped up to exactly $K$ elements by
alternating between the next-best momentum of each list, so every full round
consumes exactly $K$ true evaluations.

The two criteria answer two different questions. $d^{c}$ selects candidates that
can refine the basin the swarm already agrees on; $d^{m}$ selects candidates that
keep the second and third basins of the elite triangle alive. Both are
$O(MD)$ and neither requires a surrogate model, so no approximation error is
introduced and no matrix inversion is performed.

Equation~\eqref{eq:dm} deliberately replaces the more familiar
``farthest from the population centroid'' rule used by sparse-evaluation
PSO/CSO frameworks. In a leader-driven algorithm only three positions influence
the search, so a distant candidate that never enters the elite triangle
contributes nothing; Section~\ref{sec:ablation} shows that crowd-level diversity
scoring converts far fewer evaluations into improvements than leader-level
scoring.

\paragraph{(c) True evaluation and leader refresh.}
For every selected wolf $i \in E$ we compute $v = f(Y_i)$ and set
\begin{equation}
X^{v}_i \leftarrow Y_i, \qquad f^{v}_i \leftarrow v
\end{equation}
(the \emph{latest} semantics; an archive-style variant that keeps the best
value is tested in the ablation). The elite triangle is then re-derived by
sorting $(X^v, f^v)$ and taking its three best entries, so $\alpha,\beta,\delta$
always correspond to genuinely evaluated positions.

\subsection{Evaluation-quota scheduling}
\label{sec:q}

The split $K_c{:}K_d$ is the single decision the framework has to make. We learn
it with tabular Q-learning:
\begin{itemize}
  \item \textbf{Actions} $A = \{0.7, 0.9, 1.0\}$: the share $s = K_c/K$ of the
        quota given to consensus refinement ($s=1$ switches minority probing off).
  \item \textbf{States} $6$ states $s = (\phi-1)\cdot 2 + \sigma$, where
        $\phi \in \{\text{early},\text{mid},\text{late}\}$ is the third of the
        budget consumed and $\sigma \in \{0,1\}$ indicates whether the best
        verified fitness has stagnated for more than $\tau$ evaluations.
  \item \textbf{Decision window} the action is held for $W$ true evaluations
        and re-selected afterwards; a window is long enough that the reward
        reflects a policy rather than noise.
  \item \textbf{Reward} scale-invariant and bounded, so one Q-table serves
        problems of any magnitude:
  \begin{equation}
  r = \begin{cases}
  10\,\dfrac{f_{\text{old}} - f_{\text{new}}}{\lvert f_{\text{old}}\rvert + \lvert f_{\text{new}}\rvert + \epsilon}, & f_{\text{new}} < f_{\text{old}},\\[4pt]
  -1, & \text{otherwise},
  \end{cases}
  \end{equation}
  where $f_{\text{old}}$, $f_{\text{new}}$ are the best verified fitness at the
  start and end of the window.
  \item \textbf{Update} $Q(s,a) \leftarrow Q(s,a) + \alpha\,[\,r + \gamma \max_{a'} Q(s',a') - Q(s,a)\,]$
        with $\alpha=0.1$, $\gamma=0.9$, and an $\epsilon$-greedy policy with
        $\epsilon$ decaying from $0.10$ to $0.02$.
\end{itemize}
The Q-table is initialised optimistically in proportion to the action values,
$Q_0(s,a) = s_a$, so the agent starts from the exploitation-dominant policy that
the fixed-quota study identifies and deviates only when the reward supports it.

\subsection{Complexity}

Per round, movement and screening cost $O(MD)$ and $O(M\log M)$, and the true
evaluations cost $K\,C_f$, where $C_f$ is the cost of one objective evaluation.
The total per-round overhead $O(MD + M\log M)$ is independent of $C_f$. For a
fixed budget of $\mathrm{MaxFE}$ true evaluations the number of rounds is
$\mathrm{MaxFE}/K$ for both GWO and RASE-GWO, so the two algorithms perform the
same number of generations and the same number of true evaluations; they differ
only in how many cheap candidate updates are used to choose the $K$ positions
that are evaluated. Table~\ref{tab:complexity} summarises the comparison.

\begin{table}[t]
\centering
\caption{Per-round cost at an identical true-evaluation quota.}
\label{tab:complexity}
\begin{tabular}{lccc}
\toprule
Algorithm & Candidates & True evals/round & Overhead \\
\midrule
GWO & $N$ & $N$ & $O(ND)$ \\
Surrogate-assisted GWO & $N_s$ & $K \ll N$ & $O(N_s^3)$ \\
RASE-GWO & $M=NK$ & $K \ll M$ & $O(MD + M\log M)$ \\
\bottomrule
\end{tabular}
\end{table}

Algorithm~\ref{alg:rase} gives the complete procedure.

\begin{table*}[t]
\centering
\caption{RASE-GWO}
\label{alg:rase}
\begin{tabular}{p{0.94\textwidth}}
\toprule
\textbf{Input:} objective $f$, dimension $D$, virtual population $M$,
per-round quota $K$, budget $\mathrm{MaxFE}$, bounds $[\mathrm{lb},\mathrm{ub}]$,
initial pool $P$ \\
\textbf{Output:} best verified fitness $f^\star$ and position $X^\star$ \\
\midrule
$Y \leftarrow P(1{:}M,:)$; \; $X^{v} \leftarrow Y$; \; $f^{v} \leftarrow \infty$;
\; $\FE \leftarrow 0$; \; $T \leftarrow \mathbf{NaN}(\mathrm{MaxFE},1)$ \\
\FOR{$i=1$ \TO $K$}
  \STATE $f^{v}_i \leftarrow f(Y_i)$; \; $X^{v}_i \leftarrow Y_i$; \; $\FE \leftarrow \FE+1$;
         \; $T_{\FE} \leftarrow \min_j f^{v}_j$
\ENDFOR
$(X_\alpha,X_\beta,X_\delta) \leftarrow$ three best entries of $(X^{v},f^{v})$ \\
$Q \leftarrow Q_0$; \; $s \leftarrow$ state$(\FE)$; \; $a \sim \epsilon$-greedy$(Q,s)$;
   \; $f_{\text{win}} \leftarrow \min f^{v}$ \\
\WHILE{$\FE < \mathrm{MaxFE}$}
  \STATE $a \leftarrow 2-2\,\FE/\mathrm{MaxFE}$; move all $M$ virtual wolves by
         Eq.~\eqref{eq:move} \hfill\COMMENT{no evaluations}
  \STATE $s_{\text{share}} \leftarrow$ share of action $a$; \;
         $K_c \leftarrow \mathrm{round}(s_{\text{share}} K)$; \; $K_d \leftarrow K-K_c$
  \STATE $E_c \leftarrow$ $K_c$ smallest $d^{c}$; \;
         $E_m \leftarrow$ $K_d$ smallest $d^{m}$ \hfill\COMMENT{Eqs.~\eqref{eq:dc}--\eqref{eq:dm}}
  \STATE $E \leftarrow \mathrm{unique}(E_c \cup E_m)$; top up to $\lvert E\rvert = K$
  \FOR{$i \in E$}
    \STATE $v \leftarrow f(Y_i)$; \; $\FE \leftarrow \FE+1$; \;
           $X^{v}_i \leftarrow Y_i$; \; $f^{v}_i \leftarrow v$; \; $T_{\FE} \leftarrow \min_j f^{v}_j$
  \ENDFOR
  \STATE refresh $(X_\alpha,X_\beta,X_\delta)$ from $(X^{v},f^{v})$
  \IF{$\FE - \FE_{\text{last}} \geq W$} \hfill\COMMENT{decision window elapsed}
     \STATE $r \leftarrow$ reward of Eq.~(Q-reward); \; $s' \leftarrow$ state$(\FE)$
     \STATE $Q(s,a) \leftarrow Q(s,a) + \alpha[r + \gamma \max_{a'} Q(s',a') - Q(s,a)]$
     \STATE $a \sim \epsilon$-greedy$(Q,s')$; \; $s \leftarrow s'$; \;
            $\FE_{\text{last}} \leftarrow \FE$; \; $f_{\text{win}} \leftarrow \min f^{v}$
  \ENDIF
\ENDWHILE
\STATE \textbf{return} $f^\star = \min_j f^{v}_j$, \; $X^\star = X^{v}_{\argmin_j f^{v}_j}$
\\
\bottomrule
\end{tabular}
\end{table*}
